\documentclass[11pt]{article}
\usepackage[T1]{fontenc}
\usepackage{lmodern,microtype,amsmath,amssymb,amsthm,mathtools}
\usepackage[letterpaper,margin=1in]{geometry}
\usepackage{xurl}
\usepackage[hidelinks]{hyperref}
\hypersetup{pdftitle={Strong Chollet inequalities for all simple graph Laplacians},pdfauthor={Yongxian Zhang}}
\setlength{\emergencystretch}{3em}
\allowdisplaybreaks[1]
\newtheorem{theorem}{Theorem}[section]
\newtheorem{lemma}[theorem]{Lemma}
\newtheorem{proposition}[theorem]{Proposition}
\newtheorem{corollary}[theorem]{Corollary}
\theoremstyle{remark}\newtheorem{remark}[theorem]{Remark}
\DeclareMathOperator{\per}{per}
\DeclareMathOperator{\tr}{tr}
\newcommand{\R}{\mathbb R}
\newcommand{\PSD}{positive semidefinite}
\title{Strong Chollet inequalities\protect\\for all simple graph Laplacians}
\author{Yongxian Zhang\thanks{School of Computer Science and Engineering,
South China University of Technology, Guangzhou, China.
Email: \href{mailto:mxymmxym1@gmail.com}{mxymmxym1@gmail.com}.
ORCID: \href{https://orcid.org/0009-0000-3864-3536}{0009-0000-3864-3536}.}}
\date{October 8, 2026}
\begin{document}
\maketitle
\begin{abstract}
We prove the degree-product form of Chollet's inequality for the Laplacian
of every finite simple unweighted graph, including every principal
submatrix with degrees retained from the original graph. The result also
holds after arbitrary nonnegative diagonal additions. For a two-connected
noncycle block, a feasible point of the full matching polytope supplies a
permanent lower bound, while a cycle expansion and a trace estimate supply
an upper bound for the entrywise square. The two logarithmic coefficients
are separated by $1/60$. Cycles are treated directly, and hereditary
strength is preserved by diagonal additions and one-point sums. Block
decomposition completes the proof for arbitrary graphs. This answers the
two all-graph questions stated by Pant and Singh. A separate complete Lean
proof establishes the same graph inequalities, including empty and singular
cases; we explain its exact scope and its differences from the written proof.
\end{abstract}

\section{Introduction and statements}\label{sec:intro}
For an $n\times n$ matrix $A=(a_{ij})$, its permanent is
\[
\per A=\sum_{\pi\in\mathfrak S_n}\prod_{i=1}^n a_{i,\pi(i)}.
\]
Chollet's conjecture asks whether
$\per(A\circ B)\le\per(A)\per(B)$ for Hermitian positive-semidefinite
matrices $A,B$, where $\circ$ is the entrywise product.
The self case takes $A=B$. Pant and Singh~\cite{PantSingh} study structural
classes of real matrices and graph Laplacians. Their Section~6 asks whether
the self inequality holds for every graph Laplacian and, more strongly,
whether its right-hand side can be replaced by the permanent times the
product of vertex degrees. We prove this stronger assertion for all finite
simple unweighted graphs and for every principal submatrix.

Let $G=(V,E)$ be a finite simple undirected graph. Write $d_v=\deg_G(v)$
and let $L_G$ be its ordinary Laplacian: $(L_G)_{vv}=d_v$, and, for
$u\ne v$, $(L_G)_{uv}=-1$ on edges and $0$ otherwise.
For $S\subseteq V$, $A[S]$ denotes the principal submatrix on $S$.
The degrees in $L_G[S]$ are always the degrees in $G$, not in $G[S]$.
The permanent and diagonal product of an empty matrix are both $1$.

\begin{theorem}\label{thm:main}
For every $G$ as above and every $S\subseteq V$,
\begin{equation}\label{eq:main}
\per\bigl(L_G[S]\circ L_G[S]\bigr)
\le \per\bigl(L_G[S]\bigr)\prod_{v\in S}d_v.
\end{equation}
More generally, for any numbers $\alpha_v\ge0$ and
$A=L_G+\operatorname{diag}(\alpha_v:v\in V)$,
\begin{equation}\label{eq:diagonal-main}
\per\bigl(A[S]\circ A[S]\bigr)
\le\per\bigl(A[S]\bigr)\prod_{v\in S}(d_v+\alpha_v).
\end{equation}
Consequently $\per(L_G\circ L_G)\le\per(L_G)^2$.
\end{theorem}

There are no connectivity, minimum-degree, nonsingularity or nonemptiness
assumptions. In particular, taking a principal submatrix does not remove
the contribution of edges leaving $S$ from its diagonal.
No assertion for arbitrary weighted graph Laplacians or arbitrary
positive-semidefinite matrices is made.

The proof has three parts. On a two-connected noncycle block, diagonal
normalization turns the problem into a comparison between $\per(I+C)$
and $\per(I-B)$, where $C=B\circ B$. A trace estimate gives the upper
coefficient $19/12$ over unordered edges, and a fractional matching gives
the lower coefficient $8/5$. Their difference is $1/60$.
Cycle and edge blocks require separate arguments. Finally, the inequality
for every principal submatrix is preserved by one-point sums, allowing
the blocks of an arbitrary graph to be assembled.

We use Lieb's established block permanent inequality~\cite{Lieb} and
Edmonds' full matching-polytope theorem~\cite{Edmonds}, with all their
required conditions stated below. The diagonal and coalescence closures
are already present in Pant--Singh~\cite[Theorem 4.3 and Appendix A]{PantSingh};
we give the needed algebra in full. The written proof was previously
archived as a research note~\cite{WrittenArchive}. The present paper
organizes that proof and its verification scope; it does not present this
editorial revision as a separate solution of the underlying questions.

\section{Two established inputs and hereditary strength}\label{sec:inputs}
A real symmetric positive-semidefinite matrix $A$ is
\emph{hereditarily strong} if
\begin{equation}\label{eq:hereditary}
\per(H\circ H)\le\per(H)\prod_i h_{ii}
\quad\text{for every principal submatrix }H\text{ of }A.
\end{equation}
The adjective refers to the actual principal diagonal entries.

\begin{theorem}[Lieb's block permanent inequality]\label{thm:lieb}
Let $A$ be Hermitian positive semidefinite and let $S_1,\ldots,S_k$
partition its index set. Then
\[
\per A\ge\prod_{j=1}^k\per A[S_j].
\]
In particular $\per A\ge\prod_i a_{ii}\ge0$.
\end{theorem}
This is the non-strict inequality proved in~\cite{Lieb}; no equality
classification or unresolved permanental-dominance assertion is used.

\begin{theorem}[Edmonds' matching polytope]\label{thm:edmonds}
For a finite simple graph $H$, a vector $(x_e)_{e\in E(H)}$ is a convex
combination of incidence vectors of matchings if and only if
\begin{align*}
x_e&\ge0,\\
\sum_{e\ni v}x_e&\le1 &&(v\in V(H)),\\
\sum_{e\in E_H(U)}x_e&\le\frac{|U|-1}{2}
&& (U\subseteq V(H),\ |U|\text{ odd}).
\end{align*}
Here $E_H(U)$ denotes the edges with both ends in $U$.
\end{theorem}
In particular, for any real edge weights $w_e$, some matching $F$
satisfies $\sum_{e\in F}w_e\ge\sum_e x_e w_e$.
Indeed this follows by averaging over the convex combination.

\section{A cycle-trace permanent upper bound}\label{sec:trace}
\begin{lemma}\label{lem:upper}
Let $C$ be real symmetric, entrywise nonnegative, with zero diagonal
and all row sums at most $1/2$. Then
\begin{equation}\label{eq:upper}
\log\per(I+C)\le\frac{19}{12}\sum_{i<j}c_{ij}^2.
\end{equation}
\end{lemma}
\begin{proof}
A permutation is a family of vertex-disjoint nontrivial directed cycles,
with its remaining vertices fixed. A two-cycle is counted once; longer
cycles are counted up to cyclic rotation with orientations distinguished.
Give such a cycle $\gamma$ weight $w(\gamma)$, the product of its entries
of $C$. Dropping disjointness only adds nonnegative terms, so
\[
\per(I+C)\le\prod_{\gamma}(1+w(\gamma))
\le \exp\left(\sum_{\gamma}w(\gamma)\right).
\]
For cycles of length $k$, the total weight is at most $\tr(C^k)/k$:
each simple directed cycle has exactly $k$ rooted closed walks in the
trace, and the trace also counts nonsimple closed walks with nonnegative
weights. This statement also applies to $k=2$.

All eigenvalues of $C$ lie in $[-1/2,1/2]$, since its row sums are at
most $1/2$. For any such eigenvalue and $k\ge2$,
$\lambda^k\le2^{-(k-2)}\lambda^2$. Thus
\[
\sum_{\gamma}w(\gamma)
\le\tr(C^2)\sum_{k\ge2}\frac{2^{-(k-2)}}k.
\]
The cycle family itself is finite; extending the nonnegative upper bound
to this convergent infinite series is harmless. Moreover
\[
\sum_{k\ge2}\frac{2^{-(k-2)}}k
\le\frac12+\frac16+\frac14\sum_{k\ge4}2^{-(k-2)}
=\frac{19}{24}.
\]
Since $\tr(C^2)=2\sum_{i<j}c_{ij}^2$, this proves
\eqref{eq:upper}. Its logarithm is defined because the identity
permutation contributes $1$ to $\per(I+C)$.
\end{proof}

\section{A feasible matching in each noncycle block}\label{sec:matching}
Here a graph is two-connected if it has at least three vertices and
remains connected after deletion of any single vertex.

\begin{lemma}\label{lem:matching}
Let $G$ be two-connected and not a cycle. Set $d_v=\deg_G(v)$ and
\[
c_{uv}=\frac1{d_ud_v},\qquad x_{uv}=\frac95 c_{uv}\quad(uv\in E(G)).
\]
Then $x$ belongs to the matching polytope of $G$.
\end{lemma}
\begin{proof}
Every $d_v\ge2$. Nonnegativity holds, and the vertex constraints follow
from
\[
\sum_{u\sim v}x_{uv}=\frac9{5d_v}\sum_{u\sim v}\frac1{d_u}
\le\frac9{10}<1.
\]
We check every odd-set constraint, including the small sets.

For $U\subseteq V(G)$ with $|U|\ge2$, call $v\in U$ exceptional if
$d_v\ge3$ or $v$ has a neighbor outside $U$. There are at least two
exceptional vertices. For proper $U$, no exception contradicts
connectivity, and a unique exception $v$ would disconnect the nonempty
sets $U\setminus\{v\}$ and $V(G)\setminus U$ upon deletion of $v$.
For $U=V(G)$, all degrees equal to two would make $G$ a cycle.
If only one vertex $v$ had degree at least three, $G-v$ would be a
connected graph of maximum degree two with at least three degree-one
vertices, namely the neighbors of $v$. A connected graph of maximum
degree two is a path or a cycle and cannot have that property.

Let $d_U(v)$ count neighbors in $U$. The inequality
$2/(d_ud_v)\le1/d_u^2+1/d_v^2$ implies
\[
4\sum_{uv\in E_G(U)}c_{uv}\le2\sum_{v\in U}\frac{d_U(v)}{d_v^2}.
\]
Each deficit $1-2d_U(v)/d_v^2$ is nonnegative. It is at least $1/3$
when $d_v\ge3$, and at least $1/2$ when $d_v=2$ has an external
neighbor. The two exceptions therefore give
\begin{equation}\label{eq:deficit}
4\sum_{uv\in E_G(U)}c_{uv}\le |U|-\frac23.
\end{equation}
For odd $|U|\ge5$, it follows that
\[
\sum_{uv\in E_G(U)}x_{uv}\le\frac{9|U|-6}{20}
\le\frac{|U|-1}{2},
\]
where the last comparison already holds for $|U|\ge4$.

If $|U|=3$ and $G[U]$ has at most two edges, its $c$-sum is at most
$1/2$, giving an $x$-sum at most $9/10$. If $G[U]$ is a triangle,
at least two of its vertices have degree at least three. Otherwise the
third vertex separates the two degree-two vertices from the rest of
$G$, unless $G$ is that triangle, which was excluded. Hence
\[
\sum_{uv\in E_G(U)}c_{uv}\le\frac16+\frac16+\frac19=\frac49,
\]
and its $x$-sum is at most $4/5$. A singleton has no internal edge.
All conditions of Theorem~\ref{thm:edmonds} hold.
\end{proof}

\begin{proposition}\label{prop:block}
The Laplacian of every two-connected noncycle graph is hereditarily strong.
For any principal subset $S$, put
$X_S=\sum_{uv\in E(G[S])}(d_ud_v)^{-2}$. More precisely,
\begin{equation}\label{eq:margin}
\per(L_G[S]\circ L_G[S])
\le e^{-X_S/60}\per(L_G[S])\prod_{v\in S}d_v.
\end{equation}
Thus the inequality is strict if $G[S]$ contains an edge, and is an
equality when $G[S]$ has no edge.
\end{proposition}
\begin{proof}
The empty case is immediate. All original degrees are positive. Normalize
\[
M=\operatorname{diag}(d_v^{-1/2}:v\in S)L_G[S]
\operatorname{diag}(d_v^{-1/2}:v\in S)=I-B.
\]
The matrix $M$ is positive semidefinite. The matrix $B$ is symmetric
with zero diagonal and $b_{uv}=1/\sqrt{d_ud_v}$ on edges of $G[S]$.
For $C=B\circ B$, its edge entry is $c_{uv}=1/(d_ud_v)\le1/4$,
and its row sums are at most $1/2$. Lemma~\ref{lem:upper} gives
\begin{equation}\label{eq:upper-normalized}
\log\per(M\circ M)=\log\per(I+C)\le\frac{19}{12}X_S.
\end{equation}

Restrict the feasible vector of Lemma~\ref{lem:matching} to $G[S]$.
Restriction preserves all matching-polytope constraints. With edge
weights $\log(1+c_{uv})$, Theorem~\ref{thm:edmonds} supplies a matching
$F$ whose weight is at least
$\frac95\sum_{uv\in E(G[S])}c_{uv}\log(1+c_{uv})$.
Partition $S$ into its matched pairs and remaining singletons. Each pair
has permanent $1+b_{uv}^2=1+c_{uv}$. Theorem~\ref{thm:lieb} yields
\[
\per M\ge\prod_{uv\in F}(1+c_{uv}).
\]
For $0\le t\le1/4$,
\[
\log(1+t)\ge\frac{2t}{2+t}\ge\frac89t.
\]
The first inequality follows by differentiation: the difference has
derivative $t^2/((1+t)(2+t)^2)\ge0$ and value zero at $0$.
Consequently
\begin{equation}\label{eq:lower-normalized}
\log\per M\ge\frac85X_S.
\end{equation}
The permanent is at least $1$ by Theorem~\ref{thm:lieb}, so these
logarithms are well defined. Subtracting
\eqref{eq:lower-normalized} from \eqref{eq:upper-normalized} gives
$\per(M\circ M)\le e^{-X_S/60}\per M$.
If $Z_S=\prod_{v\in S}d_v$, diagonal scaling gives
$\per L_G[S]=Z_S\per M$ and
$\per(L_G[S]\circ L_G[S])=Z_S^2\per(M\circ M)$.
This proves \eqref{eq:margin}. An internal edge makes $X_S>0$;
without any internal edge $M=I$, giving equality.
\end{proof}

\section{Cycle and bridge blocks}\label{sec:cycles}
\begin{lemma}\label{lem:bipartite}
A real positive-semidefinite matrix with nonpositive off-diagonal entries
and bipartite off-diagonal support is hereditarily strong.
\end{lemma}
\begin{proof}
Fix a principal submatrix $H$ and a permanent monomial
$t_\pi=\prod_i h_{i,\pi(i)}$. Every nonzero nontrivial permutation cycle
in its support has even length, so $t_\pi\ge0$. Positivity of every
two-by-two principal minor gives $|h_{ij}|\le\sqrt{h_{ii}h_{jj}}$.
Therefore, with $D=\prod_i h_{ii}$,
\[
0\le t_\pi\le\prod_i\sqrt{h_{ii}h_{\pi(i),\pi(i)}}=D.
\]
Summing $t_\pi^2\le D t_\pi$ proves \eqref{eq:hereditary}.
The argument includes zero diagonal entries and the empty submatrix.
\end{proof}

\begin{lemma}\label{lem:cycles}
The Laplacians of $K_2$ and every cycle are hereditarily strong.
\end{lemma}
\begin{proof}
Lemma~\ref{lem:bipartite} handles $K_2$, even cycles and every proper
principal submatrix of a cycle Laplacian: the latter has forest support.
For the full odd cycle $C_n$, $n\ge3$, put $M=L_{C_n}/2=I-B$,
where $b_{uv}=1/2$ on cycle edges. A maximum matching has $(n-1)/2$
edges, so Theorem~\ref{thm:lieb} and the preceding scalar bound give
\[
\log\per M\ge\frac{n-1}{2}\log\frac54\ge\frac{n-1}{9}.
\]
For $C=B\circ B$, the simple directed cycles are the $n$ two-cycles
of weight $1/16$ and the two orientations of the full cycle of weight
$4^{-n}$. The cycle-product argument of Lemma~\ref{lem:upper} gives
\[
\log\per(M\circ M)\le\frac n{16}+\frac2{4^n}.
\]
For every integer $n\ge3$,
\[
\frac{n-1}{9}-\frac n{16}-\frac2{4^n}
=\frac{7n-16}{144}-\frac2{4^n}>0.
\]
The expression is increasing in $n$ and equals $1/288$ at $n=3$.
Undoing normalization proves the full-cycle inequality, strictly.
\end{proof}

\section{Diagonal additions and one-point sums}\label{sec:closure}
The following closure arguments also account for all principal subsets
at articulation vertices, where degrees from both sides contribute.

\begin{lemma}\label{lem:diagonal}
Hereditary strength is preserved by adding a nonnegative diagonal matrix.
\end{lemma}
\begin{proof}
It suffices to add $\alpha\ge0$ at one diagonal index $v$. For any
principal submatrix containing $v$, write its old matrix as $A$ and put
\[
a=a_{vv},\quad P=\per A,\quad R=\per A(v),\quad
Q=\per(A\circ A),\quad T=\per(A(v)\circ A(v)),
\]
where $A(v)$ deletes $v$, and let $h$ be its diagonal product.
Expansion in the changed diagonal entry gives
\[
\per(A+\alpha E_{vv})=P+\alpha R,\qquad
\per((A+\alpha E_{vv})\circ(A+\alpha E_{vv}))
=Q+(2a\alpha+\alpha^2)T.
\]
Hereditary strength gives $Q\le ahP$ and $T\le hR$, while
Theorem~\ref{thm:lieb} gives $P\ge aR$. All quantities multiplied
below are nonnegative. Hence
\begin{align*}
Q+(2a\alpha+\alpha^2)T
&\le h\bigl[aP+(2a\alpha+\alpha^2)R\bigr]\\
&\le h(a+\alpha)(P+\alpha R),
\end{align*}
because the second difference is $h\alpha(P-aR)\ge0$.
Other principal subsets are unchanged. Iteration proves the lemma.
\end{proof}

\begin{lemma}\label{lem:onepoint}
The one-point sum of two hereditarily strong positive-semidefinite
matrices is hereditarily strong. Direct sums have the same property.
\end{lemma}
\begin{proof}
For the one-point sum, identify distinguished indices of $A_1,A_2$,
add their diagonal entries $a_1,a_2$, and set cross entries between
the two remaining index sets to zero. The resulting matrix $A$ is
positive semidefinite, being the sum of the two embedded matrices.
Let $B_i$ delete the distinguished index and put
\[
P_i=\per A_i,\quad R_i=\per B_i,\quad
Q_i=\per(A_i\circ A_i),\quad T_i=\per(B_i\circ B_i),
\]
with $h_i$ the diagonal product of $B_i$.
A permutation cycle through the common vertex cannot visit both sides:
such a cycle would have to pass through that vertex twice. Expansion
by this cycle, including the fixed-vertex case, gives
\begin{align}
\per A&=P_1R_2+R_1P_2,\label{eq:onepoint-per}\\
\per(A\circ A)&=Q_1T_2+T_1Q_2+2a_1a_2T_1T_2.
\label{eq:onepoint-square}
\end{align}
The additional term in \eqref{eq:onepoint-square} comes from
$(a_1+a_2)^2-a_1^2-a_2^2$. Hereditary strength and nonnegativity imply
\[
\per(A\circ A)\le h_1h_2
\bigl(a_1P_1R_2+a_2R_1P_2+2a_1a_2R_1R_2\bigr).
\]
The target right-hand side is
$h_1h_2(a_1+a_2)(P_1R_2+R_1P_2)$. After removing $h_1h_2$, their
difference is
\[
a_2R_2(P_1-a_1R_1)+a_1R_1(P_2-a_2R_2)\ge0
\]
by Theorem~\ref{thm:lieb}. This proves the whole-matrix inequality.

For a principal subset containing the common vertex, take its selected
indices from each side together with that vertex on each side, and apply
the same computation. Either selected side may consist just of the common
vertex; the empty-matrix conventions give $R_i=T_i=h_i=1$ there.
For a subset omitting the common vertex, the matrix is a direct sum.
Permanents and diagonal products multiply on direct sums, proving both
remaining assertions, with empty and singular cases included.
\end{proof}

\section{Assembly of all graphs}\label{sec:assembly}
\begin{proof}[Proof of Theorem~\ref{thm:main}]
The standard block decomposition of a connected graph expresses it as
successive one-point sums of its bridges and maximal two-connected
blocks; see, for example,~\cite[Chapter 3]{Diestel}.
At each articulation vertex, its full degree is the sum of its degrees
in those blocks. Thus the same decomposition identifies the actual
Laplacian, with its full diagonal, as their matrix one-point sum.
Each bridge is $K_2$. Each two-connected block is a cycle or a noncycle.
Proposition~\ref{prop:block} and Lemma~\ref{lem:cycles} establish
hereditary strength for all blocks, and Lemma~\ref{lem:onepoint}
preserves it during assembly. A one-vertex component has zero Laplacian:
its nonempty permanent is zero, and its empty permanent is one.
Disconnected graphs are direct sums of their connected components.
This proves \eqref{eq:main} for all $G,S$.

Lemma~\ref{lem:diagonal} proves \eqref{eq:diagonal-main}. Finally
Theorem~\ref{thm:lieb} gives $\per L_G\ge\prod_v d_v\ge0$.
Multiplication by the nonnegative $\per L_G$ yields
$\per(L_G\circ L_G)\le\per(L_G)^2$.
\end{proof}

\begin{remark}
Proposition~\ref{prop:block} records the quantitative gap already
implicit in the block comparison. It classifies equality for those
principal submatrices of a two-connected noncycle graph. It is not a
classification of equality for all assembled graphs, and the constant
$1/60$ is not claimed to be optimal.
\end{remark}

\section{Formal verification and reproduction}\label{sec:formal}
The same graph endpoint has a complete independent Lean proof in the
fixed package~\cite{LeanArchive}. Its main declaration is
\begin{quote}\small\path{Chollet.strongChollet_all_graphs}\end{quote}
in \path{src/CholletAllGraphs.lean}. It uses the actual Mathlib
\path{SimpleGraph.lapMatrix} and \path{Matrix.permanent}, and the shared
\path{StrongChollet} definition quantifies over every principal vertex
subset, retaining the degrees of the original graph. The stronger
declaration \path{laplacianDiagonalStrong_all} permits arbitrary
nonnegative diagonal additions to the whole graph Laplacian.

Lean is pinned to 4.34.1, compiler commit
\nolinkurl{5045d0056413266e57c625dcd7c365b10e377c52}, and Mathlib to
\nolinkurl{d13f23b723b8a846827a245b89c10fc7d3f11612}.
The fixed records report fresh compilation of all 83 source modules and
empty-kernel, trust-level-zero replay of 985 owned declarations and
their 44,173-declaration transitive closure. Only the standard axioms
\path{propext}, \path{Classical.choice} and \path{Quot.sound} are
allowed, and a deliberately false proof is rejected.
These are records of the frozen verification run, not a claim that this
editorial revision reran the kernel.

The Lean proof is not a line-by-line transcription of this paper.
It derives the matching lower bound from proved Gram and Fischer
identities, uses the bundled proved perfect-matching polytope theorem,
and bounds degree-two vertex-robust blocks through a vanishing third
trace and scaled feasible weights. Vertex-count induction strengthened
by nonnegative diagonal additions then replaces explicit block-tree
assembly. A semantic review checks that the unconditional endpoint
has exactly the graph quantifiers in Theorem~\ref{thm:main}.
Neither the exponential refinement in Proposition~\ref{prop:block}
nor every intermediate written argument is claimed to be separately
formalized. Third-party source adaptations and licenses are documented
in the package's \path{THIRD_PARTY.md}.

For reproduction, start with the package README and \path{case.json}.
The latter fixes all source hashes and the complete module order.
Run \path{lake build} for the candidate build. The included
\path{reproduce.py} documents fresh compilation and trust-zero replay
with the pinned official dependencies and a new output directory.
Its recorded results are \path{verification/final-review.json},
\path{verification/replay-summary.json} and
\path{verification/negative-kernel.log}. The traditional proof above
does not use a finite computation as a premise.

\section{Further questions}
The graph theorem leaves the general positive-semidefinite Chollet
conjecture untouched. An extension to arbitrary positive edge weights
would require additional control: the present fractional matching uses
unweighted vertex degrees and the fixed normalized edge size. The proof
also does not provide a full equality classification after block assembly.
Those questions are distinct from the two all-graph unweighted questions
answered here.

\section*{Funding}
This research was conducted independently and received no external funding.
\section*{AI-assisted research methods}
OpenAI AI assistants assisted mathematical research, proof development,
Lean formalization and internal review. Claims are supported by the
written proof and the fixed formal-verification evidence. AI output is
not itself proof evidence, and internal review is not external human
peer review.
\section*{Declaration of generative AI and AI-assisted technologies
in the writing process}
OpenAI AI assistants assisted drafting, organization and LaTeX editing.
The author takes responsibility for the manuscript and its account of
the verification evidence.

\begin{thebibliography}{9}
\raggedright
\bibitem{PantSingh}
P. Pant and R. Singh, \emph{Structural Classes for Chollet's Permanent
Conjecture}, arXiv:2604.24192v2 (2026), revised 18 September 2026,
Section 6, Theorem 4.3 and Appendix A.
\url{https://arxiv.org/abs/2604.24192v2}.
\bibitem{Lieb}
E. H. Lieb, \emph{Proofs of some conjectures on permanents},
Journal of Mathematics and Mechanics \textbf{16} (1966), 127--134.
\url{https://doi.org/10.1512/iumj.1967.16.16008}.
\bibitem{Edmonds}
J. Edmonds, \emph{Maximum matching and a polyhedron with 0,1-vertices},
Journal of Research of the National Bureau of Standards B \textbf{69B}
(1965), 125--130. \url{https://doi.org/10.6028/jres.069B.013}.
\bibitem{Diestel}
R. Diestel, \emph{Graph Theory}, fifth edition, Graduate Texts in
Mathematics 173, Springer, 2017.
\bibitem{WrittenArchive}
Y. Zhang, \emph{A strong Chollet inequality for every simple graph
Laplacian}, written proof archive, version 1 (2026).
\url{https://doi.org/10.5281/zenodo.23248359}.
\bibitem{LeanArchive}
Y. Zhang, \emph{Strong Chollet inequalities for all finite simple graph
Laplacians: complete Lean proof}, proof archive, version 2.0 (2026).
\url{https://doi.org/10.5281/zenodo.23252138}.
Fixed source commit \nolinkurl{4d2eefd40ee930216ccd8fc0f51e4bf694251967}.
\end{thebibliography}
\end{document}
